%% kapa: exam=1 round=27 year=2016 period=2 form=A subject=appraisal_law subject_ko=감정평가관계법규 seq=19 number=19
%% source: Q-Net 감정평가사 기출문제 (cst00309) artlSeq=5169402 file=extracted/1008/5169402/제27회 감정평가사 1차 문제/제27회 감정평가사 2교시문제(A형).pdf
%% conversion: kordoc 4.24.51 -> tools/gen1.py (auto); review=proofread-sonnet; flags=none
\begin{question}{19}
부동산가격공시 및 감정평가에 관한 법령상 시장․군수가 개별공시지가를 정정할 수 있는 사유가 아닌 것은?
\begin{choices}
  \item 표준지 선정의 착오
  \item 개별공시지가를 결정․공시하기 위하여 개별토지가격을 산정한 때에 토지소유자의 의견청취절차를 거치지 아니한 경우
  \item 토지가격비준표의 적용에 오류가 있는 경우
  \item 용도지역 등 토지가격에 영향을 미치는 주요 요인의 조사를 잘못한 경우
  \item 토지가격이 전년대비 급격하게 상승한 경우
\end{choices}
\end{question}
\answer{5}
